% =====================================================================
% Notas de desenvolvimento do TP — base para o relatório SBC
% Cada pessoa edita SÓ o próprio arquivo em secoes/ (evita conflito no Git).
% Compilar: pdflatex notas && bibtex notas && pdflatex notas && pdflatex notas
% (ou faça upload da pasta relatorio/notas/ no Overleaf)
% =====================================================================
\documentclass[11pt,a4paper]{article}

\usepackage[utf8]{inputenc}
\usepackage[T1]{fontenc}
\usepackage[brazil]{babel}
\usepackage[margin=2.2cm]{geometry}
\usepackage{amsmath,amssymb}
\usepackage{graphicx}
\usepackage{booktabs,tabularx,array}
\newcolumntype{L}[1]{>{\raggedright\arraybackslash}p{#1}}
\newcolumntype{Y}{>{\raggedright\arraybackslash}X}
\usepackage{enumitem}
\usepackage{xcolor}
\usepackage[hidelinks]{hyperref}

\graphicspath{{figuras/}}

% ---------------------------------------------------------------------
% Comandos do grupo (use em todas as seções)
% ---------------------------------------------------------------------
% Código, nomes de função e arquivos: \codigo{bfs}
\newcommand{\codigo}[1]{\texttt{#1}}
% Situação de cada item: \feito, \andamento, \afazer
\newcommand{\feito}{\textcolor{green!50!black}{\textbf{feito}}}
\newcommand{\andamento}{\textcolor{orange!80!black}{\textbf{em andamento}}}
\newcommand{\afazer}{\textcolor{gray}{a fazer}}
% Entrada no registro: \registro{AAAA-MM-DD}{o que foi feito; commit ou PR}
\newcommand{\registro}[2]{\item[\textbf{#1}] #2}
% Valor ainda não obtido, na coluna "Obtido" das tabelas: \falta
\newcommand{\falta}{\textcolor{red!70!black}{\textbf{?}}}
% Lembrete visível no PDF: \pendente{falta medir o tempo}
\newcommand{\pendente}[1]{\textcolor{red!70!black}{[PENDENTE: #1]}}

\title{TP Grafos --- Análise de Redes Sociais\\[2pt]
       \large Notas de desenvolvimento (base para o relatório)}
\author{Brunno, Vitória, Iago, Arthur, Ravi e Kaio}
\date{Atualizado em \today}

\begin{document}
\maketitle

\section*{Como usar este documento}

\begin{itemize}[leftmargin=*]
  \item Cada pessoa escreve \textbf{só no próprio arquivo} em \codigo{secoes/}. Assim ninguém
        edita o mesmo arquivo ao mesmo tempo e o Git não gera conflito.
  \item Atualize as notas \textbf{no mesmo PR do código}: implementou uma função, registre nas
        notas antes de abrir o PR. Commit sugerido: \codigo{docs: atualiza notas da frente N}.
  \item Escreva em frases completas, como no relatório. O texto daqui vai ser reaproveitado quase
        sem mudança na seção de cada um.
  \item Toda referência entra em \codigo{referencias.bib} e é citada com \codigo{\textbackslash cite\{chave\}}.
        Figuras (gráficos, prints, diagramas) vão em \codigo{figuras/}.
  \item Para valores numéricos, registre sempre \emph{obtido} e \emph{referência}, com a base e os parâmetros.
        Os caminhos nas tabelas são relativos à pasta \codigo{redes/}.
\end{itemize}

\paragraph{Problema aplicado.} Identificar os influenciadores da rede da Twitch em português e
verificar qual centralidade os encontra melhor, comparando cada ranking, pelo $\tau$ de Kendall,
com o alcance numa simulação SIR e com as visualizações reais dos streamers~\cite{kitsak2010}.

\paragraph{Restrição do enunciado.} É proibido usar bibliotecas prontas para criar ou processar o
grafo, como o NetworkX, salvo para comparação de desempenho. No projeto, o NetworkX só aparece em
\codigo{desempenho/}.

% =====================================================================
% Frente 1 --- Estrutura e dados (Brunno)
% =====================================================================
\section{Frente 1 --- Estrutura do grafo e dados}

\noindent\textbf{Responsável:} Brunno.
\textbf{No relatório:} introdução; bases e modelagem; responsabilidades; revisão final.
\textbf{Na Etapa 1:} justificativa da modelagem, diagrama de classes e montagem do documento.

\subsection{O que foi implementado}

{\small\begin{tabularx}{\linewidth}{@{}L{3.9cm} L{4.7cm} Y L{1.9cm}@{}}
\toprule
Função / classe & Arquivo & Complexidade & Situação \\
\midrule
\codigo{Grafo} & \codigo{nucleo/grafo.py} & inserir nó, inserir ou remover aresta e testar aresta: $O(1)$ médio; \codigo{arestas()}: $O(n+m)$; memória $O(n+m)$ & \feito \\
\codigo{IndexadorDeNos} & \codigo{nucleo/indexador.py} & $O(n)$ & \feito \\
\codigo{para\_listas\_indexadas} & \codigo{nucleo/indexador.py} & $O(n+m)$ & \feito \\
\codigo{carregar\_karate}, \codigo{carregar\_clubes\_karate} & \codigo{leitores/bases.py} & $O(m)$ & \feito \\
\codigo{carregar\_twitch\_ptbr} & \codigo{leitores/bases.py} & $O(n+m)$ & \feito \\
\codigo{carregar\_facebook}, \codigo{carregar\_thrones} & \codigo{leitores/bases.py} & $O(m)$ & \feito \\
\codigo{erdos\_renyi} & \codigo{geradores/modelos.py} & $O(n+m)$~\cite{batagelj2005} & \feito \\
\codigo{barabasi\_albert} & \codigo{geradores/modelos.py} & $O(n \cdot m)$ esperado~\cite{barabasi1999} & \feito \\
\codigo{watts\_strogatz} & \codigo{geradores/modelos.py} & $O(n \cdot k)$ esperado~\cite{watts1998} & \feito \\
\bottomrule
\end{tabularx}}

\subsection{Decisões e justificativas}

\begin{itemize}[leftmargin=*]
  \item \textbf{Lista de adjacência, e não matriz.} As redes sociais são esparsas: na Twitch PT,
        a densidade é $2m/(n(n-1)) \approx 0{,}017$, ou seja, uma matriz teria 98\% de zeros.
        A lista ocupa $O(n+m)$ de memória, contra $O(n^2)$ da matriz (3,6 milhões de células só na
        Twitch PT), e a BFS custa $\Theta(n+m)$ com lista e $\Theta(n^2)$ com
        matriz~\cite{cormen2022,sedgewick2011}.
  \item \textbf{Dicionário de dicionários} (\codigo{\_adj[v][w] = peso}). Testar se uma aresta existe
        custa $O(1)$ em média~\cite{pythoncomplexidade}, e o peso fica junto do vizinho. O peso é
        necessário na Network of Thrones e no Louvain, que cria laços com peso ao agregar
        comunidades~\cite{blondel2008}. O NetworkX usa a mesma estrutura pelos mesmos
        motivos~\cite{hagberg2008}.
  \item \textbf{Predecessores (\codigo{\_pred}) no grafo direcionado}, para o grau de entrada e o
        PageRank, que soma sobre as arestas que chegam em cada nó~\cite{brin1998}. No não
        direcionado, \codigo{\_pred} é o próprio \codigo{\_adj}, sem cópia.
  \item \textbf{Contador de arestas} \codigo{\_m}: \codigo{numero\_de\_arestas()} em $O(1)$, usado na
        densidade, no grau médio e na modularidade.
  \item \textbf{Laços:} aceitos na classe (o Louvain precisa) e descartados pelos leitores. O laço
        conta 2 no grau, para a soma dos graus ser sempre $2m$; o NetworkX 3.7 segue a mesma convenção.
  \item \textbf{\codigo{\_\_slots\_\_}} para economizar memória por objeto; \codigo{vizinhos()} devolve
        um iterador e \codigo{arestas()} é um gerador, para não copiar listas.
  \item \textbf{Reindexação $0..n-1$:} algoritmos pesados (Brandes, PageRank) usam listas indexadas
        por inteiro, de menor custo constante que dicionários. \pendente{medir dict vs. lista indexada
        numa BFS para sustentar essa afirmação}
  \item \textbf{Nomes:} vértices \codigo{v} (origem) e \codigo{w} (destino), como em~\cite{sedgewick2011};
        código e pastas em português.
  \item \textbf{Remoção de aresta:} \codigo{remover\_aresta(v, w)} confere se a aresta existe antes de
        apagar (senão lança \codigo{KeyError} e o grafo fica intacto), apaga os dois sentidos e
        decrementa o contador. Os nós continuam no grafo. Quem usa é o Watts--Strogatz.
  \item \textbf{Leitores:} os caminhos padrão partem da raiz do repositório (\codigo{RAIZ}), para
        funcionar de qualquer pasta; todos descartam laços. Na Twitch, os atributos se ligam aos nós
        pela coluna \codigo{new\_id} (a coluna \codigo{id} é o ID original e não aparece nas arestas).
        O CSV da Network of Thrones separa as linhas só com \codigo{\textbackslash r}, por isso é
        aberto com \codigo{newline=""}. Karate e ego-Facebook têm o mesmo formato (\codigo{v w} por
        linha) e usam a mesma função auxiliar.
  \item \textbf{Geradores:} todos recebem \codigo{semente} e usam \codigo{random.Random(semente)}, nunca o
        \codigo{random} global, para o grafo se repetir. O Erdős--Rényi pula direto para o próximo par
        sorteado (pulo geométrico), em $O(n+m)$ em vez de $O(n^2)$~\cite{batagelj2005}. O
        Barabási--Albert sorteia numa lista em que cada nó aparece uma vez por aresta, o que já dá
        chance proporcional ao grau; o total é $m(n-m)$ arestas. O Watts--Strogatz monta o anel e
        religa cada aresta com probabilidade $p$, sem criar laço nem aresta repetida; o total
        continua $n k / 2$.
\end{itemize}

\subsection{Corretude: testes e resultados}

{\small\begin{tabularx}{\linewidth}{@{}Y L{2.5cm} L{2.2cm} L{3.8cm}@{}}
\toprule
Teste & Base & Obtido & Referência \\
\midrule
Grafo completo $K_5$: arestas e grau de cada nó & sintético & 10 e 4 & $n(n-1)/2 = 10$ e $n-1 = 4$ \\
Soma dos graus $= 2m$ (Hypothesis, grafos aleatórios com laços) & sintético & passou & lema do aperto de mãos \\
Entradas $=$ saídas $= m$ no direcionado (Hypothesis) & sintético & passou & --- \\
Nós e arestas & Karate & 34 e 78 & 34 e 78~\cite{zachary1977} \\
Nós e arestas & Twitch PT & 1.912 e 31.299 & 1.912 e 31.299~\cite{rozemberczki2021} \\
Parceiros (\codigo{partner}) & Twitch PT & 279 & 279 \\
Nós e arestas & ego-Facebook & 4.039 e 88.234 & 4.039 e 88.234~\cite{mcauley2012} \\
Nós e arestas; grau do Tyrion & Network of Thrones (livro 3) & 303 e 1.008; 59 & 303 e 1.008 (arquivo usado) \\
Mesmas arestas que o NetworkX lendo o mesmo arquivo & as 4 bases & iguais & --- \\
ER com $p = 0$ e $p = 1$ & sintético & 0 e $n(n-1)/2$ & vazio e completo \\
ER com $n = 1.000$, $p = 0{,}01$, \codigo{semente=1} & sintético & 4.963 & 4.995 esperado ($\pm 5\%$) \\
BA com $n = 1.000$, $m = 3$ & sintético & 2.991, conexo & $m(n - m) = 2.991$ \\
WS com $n = 1.000$, $k = 4$ e $p \in \{0{,}1;\ 0{,}5;\ 1\}$ & sintético & 2.000 & $n k / 2 = 2.000$ \\
Mesma \codigo{semente} repete o grafo (ER, BA e WS) & sintético & passou & --- \\
Remover todas as arestas deixa $m = 0$ e os nós (Hypothesis) & sintético & passou & --- \\
\bottomrule
\end{tabularx}}

\medskip
\noindent Observação para o relatório: o artigo da Network of Thrones~\cite{beveridge2016} cita
107 personagens e 353 arestas no livro 3; o arquivo do repositório dos autores tem 303 nós e
1.008 arestas. Citar os números do arquivo usado.

\subsection{Figuras e tabelas para o relatório}
\begin{itemize}[leftmargin=*]
  \item Diagramas de classes (núcleo) e de pacotes, em \codigo{documentos/uml/} (PlantUML e PNG). \feito
  \item Tabela das 4 bases (nós, arestas, densidade e papel), no documento da Etapa 1. \feito
\end{itemize}

\subsection{Limitações e pendências}
\begin{itemize}[leftmargin=*]
  \item O Barabási--Albert começa com $m$ nós isolados (como no modelo original); o primeiro nó novo
        se liga a todos eles, e o grafo final é conexo.
  \item Exportar o grafo e os resultados em JSON para a interface, se o Kaio precisar. \afazer
\end{itemize}

\subsection{Registro}
\begin{description}[leftmargin=2.6cm,style=nextline]
  \registro{2026-11-06}{Classe \codigo{Grafo}, reindexação e 13 testes (PR \#2).}
  \registro{2026-11-09}{Nomes do código e das pastas traduzidos para português (PR \#3).}
  \registro{2026-11-10}{Classe reescrita e comentada; vértices padronizados como \codigo{v} e \codigo{w} (PR \#4). 21 testes passando.}
  \registro{2026-11-11}{Notas de desenvolvimento em LaTeX para o grupo (este documento).}
  \registro{2026-11-12}{Leitores das quatro bases e 8 testes, com os tamanhos conferidos nos arquivos reais (branch \codigo{feat/leitores}).}
  \registro{2026-11-12}{\codigo{remover\_aresta} na classe \codigo{Grafo}; geradores ER, BA e WS; 20 testes novos (branch \codigo{feat/geradores}). 49 testes passando.}
  \registro{2026-11-12}{Diagramas de classes e de pacotes; documento da Etapa 1 com bases, modelagem e responsabilidades (branch \codigo{docs/etapa1}).}
\end{description}

% =====================================================================
% Frente 2 --- Caminhos e proximidade (Vitória)
% =====================================================================
\section{Frente 2 --- Grau, caminhos e proximidade}

\noindent\textbf{Responsável:} Vitória.
\textbf{No relatório:} trabalhos relacionados; caminhos e proximidade.
\textbf{Na Etapa 1:} plano de desenvolvimento com os marcos.

\subsection{O que foi implementado}

{\small\begin{tabularx}{\linewidth}{@{}L{3.9cm} L{4.7cm} Y L{1.9cm}@{}}
\toprule
Função & Arquivo & Complexidade & Situação \\
\midrule
\codigo{centralidade\_de\_grau}, \codigo{distribuicao\_de\_graus} & \codigo{metricas/grau.py} & $O(n)$ & \andamento \\
\codigo{densidade}, \codigo{grau\_medio} & \codigo{metricas/grau.py} & $O(1)$ & \andamento \\
\codigo{bfs}, \codigo{ordem\_dfs} & \codigo{metricas/caminhos.py} & $O(n+m)$ & \afazer \\
\codigo{componentes\_conexas} & \codigo{metricas/caminhos.py} & $O(n+m)$ & \afazer \\
\codigo{dijkstra} & \codigo{metricas/caminhos.py} & $O((n+m)\log n)$~\cite{dijkstra1959} & \afazer \\
\codigo{distancia\_media}, \codigo{diametro} & \codigo{metricas/caminhos.py} & $O(n(n+m))$ & \afazer \\
\codigo{proximidade}, \codigo{harmonica} & \codigo{metricas/proximidade.py} & $O(n(n+m))$~\cite{boldi2014} & \afazer \\
\bottomrule
\end{tabularx}}

\subsection{Decisões e justificativas}
\begin{itemize}[leftmargin=*]
  \item \pendente{ex.: BFS com \codigo{collections.deque}, porque \codigo{list.pop(0)} custa $O(n)$}
  \item \pendente{ex.: DFS iterativa, porque o limite de recursão do Python estoura em caminhos longos}
\end{itemize}

\subsection{Corretude: testes e resultados}

{\small\begin{tabularx}{\linewidth}{@{}Y L{2.5cm} L{2.2cm} L{3.8cm}@{}}
\toprule
Teste & Base & Obtido & Referência \\
\midrule
Densidade, grau médio e grau máximo & Karate & \falta & 0,139; 4,588; 17 \\
Distância média e diâmetro & Karate & \falta & 2,4082 e 5 \\
Closeness dos nós 0, 2 e 33 & Karate & \falta & 0,569; 0,5593; 0,55 \\
Harmônica dos nós 33 e 0 & Karate & \falta & 23,25 e 23,167 \\
Componentes, distância média e diâmetro & Twitch PT & \falta & 1; 2,5324; 7 \\
Closeness máxima & Twitch PT & \falta & nó 127 (0,6009) \\
\bottomrule
\end{tabularx}}

\subsection{Figuras e tabelas para o relatório}
\begin{itemize}[leftmargin=*]
  \item \pendente{ex.: histograma da distribuição de graus da Twitch PT}
\end{itemize}

\subsection{Limitações e pendências}
\begin{itemize}[leftmargin=*]
  \item \pendente{ex.: distância média exata inviável no ego-Facebook; usar amostra de origens}
\end{itemize}

\subsection{Registro}
\begin{description}[leftmargin=2.6cm,style=nextline]
  \registro{2026-11-11}{Métricas de grau e 13 testes (PR \#6, em revisão).}
\end{description}

% =====================================================================
% Frente 3 --- Centralidades e desempenho (Iago)
% =====================================================================
\section{Frente 3 --- Centralidades e desempenho}

\noindent\textbf{Responsável:} Iago.
\textbf{No relatório:} centralidades; avaliação de desempenho.
\textbf{Na Etapa 1:} diagrama de sequência.

\subsection{O que foi implementado}

{\small\begin{tabularx}{\linewidth}{@{}L{3.9cm} L{4.7cm} Y L{1.9cm}@{}}
\toprule
Função & Arquivo & Complexidade & Situação \\
\midrule
\codigo{intermediacao} (exata) & \codigo{metricas/centralidade.py} & $O(nm)$~\cite{brandes2001} & \feito \\
\codigo{intermediacao} ($k$ fontes) & \codigo{metricas/centralidade.py} & $O(km)$~\cite{brandes2007} & \feito \\
\codigo{pagerank} & \codigo{metricas/centralidade.py} & $O(m)$ por iteração~\cite{brin1998} & \feito \\
\codigo{numero\_de\_nucleo}, \codigo{k\_nucleo} & \codigo{metricas/nucleo\_k.py} & $O(m)$~\cite{batagelj2003} & \feito \\
\codigo{medir.py} (tempo, memória, TEPS) & \codigo{desempenho/} & --- & \afazer \\
\codigo{comparar\_networkx.py} & \codigo{desempenho/} & --- & \afazer \\
\bottomrule
\end{tabularx}}

\subsection{Decisões e justificativas}
\begin{itemize}[leftmargin=*]
  \item Listas indexadas $0..n-1$ (\codigo{IndexadorDeNos}) no Brandes, no PageRank e no k-core: mais rápidas que os dicionários do \codigo{Grafo}.
  \item No PageRank, nós sem saída repartem o valor entre todos os nós, e a aresta é escolhida em proporção ao peso.
  \item Se o PageRank não converge em \codigo{max\_iteracoes} (acontece em grafos bipartidos, como a estrela, com a tolerância $10^{-8}$), a função emite um aviso e devolve a última iteração, com erro da ordem de $0{,}85^{100} \approx 10^{-7}$, em vez de lançar exceção.
  \item No k-core, laços são ignorados; no direcionado, o grau é entrada + saída, como no NetworkX.
\end{itemize}

\subsection{Corretude: testes e resultados}

{\small\begin{tabularx}{\linewidth}{@{}Y L{2.5cm} L{2.2cm} L{3.8cm}@{}}
\toprule
Teste & Base & Obtido & Referência \\
\midrule
Betweenness normalizada dos nós 0, 33 e 32 & Karate & 0,4376; 0,3041; 0,1452 & 0,4376; 0,3041; 0,1452 \\
PageRank ($d = 0{,}85$) dos nós 33, 0 e 32; soma & Karate & 0,1009; 0,097; 0,0717; 1 & 0,1009; 0,097; 0,0717; 1 \\
Núcleo máximo & Karate & 4 (os mesmos 10 nós) & 4 (nós 0, 1, 2, 3, 7, 8, 13, 30, 32, 33) \\
Betweenness e PageRank máximos & Twitch PT & \falta & nó 127 (0,0993 e 0,0119) \\
Núcleo máximo & Twitch PT & \falta & 37, com 154 nós \\
Centro da estrela; grafo completo & sintético & $(n-1)(n-2)/2$; 0 & $(n-1)(n-2)/2$; 0 \\
\bottomrule
\end{tabularx}}

\subsection{Desempenho}
\begin{itemize}[leftmargin=*]
  \item Ambiente: \pendente{CPU, RAM, sistema, versão do Python e do NetworkX}
  \item Metodologia: mediana de 5 execuções com \codigo{time.perf\_counter}; pico de memória com
        \codigo{tracemalloc}; grafos ER, BA e WS de $10^3$ a $10^5$ nós e as 4 bases.
  \item Resultados: \pendente{tabela nosso código vs. NetworkX; inclinação do log-log}
\end{itemize}

\subsection{Figuras e tabelas para o relatório}
\begin{itemize}[leftmargin=*]
  \item \pendente{gráfico log-log de tempo $\times$ $(n+m)$, com retas $O(n+m)$ e $O(nm)$}
\end{itemize}

\subsection{Limitações e pendências}
\begin{itemize}[leftmargin=*]
  \item \pendente{ex.: betweenness exata inviável em grafos grandes; usar amostragem~\cite{brandes2007}}
\end{itemize}

\subsection{Registro}
\begin{description}[leftmargin=2.6cm,style=nextline]
  \registro{2026-11-13}{Diagrama de sequência da Etapa 1 e intermediação de Brandes (exata e por $k$ origens), PageRank e k-core, com testes.}
\end{description}

% =====================================================================
% Frente 4 --- Agrupamento e comunidades (Arthur)
% =====================================================================
\section{Frente 4 --- Agrupamento e comunidades}

\noindent\textbf{Responsável:} Arthur.
\textbf{No relatório:} comunidades; limitações e ética.
\textbf{Na Etapa 1:} diagrama de atividades do Louvain.

\subsection{O que foi implementado}

{\small\begin{tabularx}{\linewidth}{@{}L{3.9cm} L{4.7cm} Y L{1.9cm}@{}}
\toprule
Função & Arquivo & Complexidade & Situação \\
\midrule
\codigo{triangulos} & \codigo{metricas/agrupamento.py} & $O(m^{1,5})$~\cite{schank2005} & \afazer \\
\codigo{agrupamento\_local}, \codigo{agrupamento\_medio} & \codigo{metricas/agrupamento.py} & \falta & \afazer \\
\codigo{transitividade} & \codigo{metricas/agrupamento.py} & \falta & \afazer \\
\codigo{modularidade} & \codigo{comunidades/} & $O(m)$~\cite{newman2004} & \afazer \\
\codigo{louvain} & \codigo{comunidades/} & $\approx O(m)$ por passada~\cite{blondel2008} & \afazer \\
\codigo{propagacao\_de\_rotulos} & \codigo{comunidades/} & $O(m)$ por iteração~\cite{raghavan2007} & \afazer \\
\codigo{nmi}, \codigo{ari} & \codigo{comunidades/} & \pendente{} ~\cite{danon2005,hubert1985} & \afazer \\
\bottomrule
\end{tabularx}}

\subsection{Decisões e justificativas}
\begin{itemize}[leftmargin=*]
  \item \pendente{ex.: fórmula do ganho de modularidade usada no Louvain}
\end{itemize}

\subsection{Corretude: testes e resultados}

{\small\begin{tabularx}{\linewidth}{@{}Y L{2.5cm} L{2.2cm} L{3.8cm}@{}}
\toprule
Teste & Base & Obtido & Referência \\
\midrule
Clustering médio, transitividade e triângulos & Karate & \falta & 0,5706; 0,2557; 45 \\
Modularidade do Louvain & Karate & \falta & entre 0,39 e 0,42 (máximo 0,4198) \\
Modularidade da divisão real (gabarito) & Karate & \falta & 0,3582 \\
Clustering médio, transitividade e triângulos & Twitch PT & \falta & 0,3199; 0,131; 173.510 \\
Louvain & Twitch PT & \falta & 0,286 a 0,296 ($\approx$ 7 comunidades) \\
Clustering médio, transitividade e triângulos & ego-Facebook & \falta & 0,6055; 0,5192; 1.612.010 \\
NMI e ARI de partições iguais & sintético & \falta & 1 e 1 \\
\bottomrule
\end{tabularx}}

\subsection{Figuras e tabelas para o relatório}
\begin{itemize}[leftmargin=*]
  \item \pendente{ex.: diagrama de atividades do Louvain; Karate colorido por comunidade}
\end{itemize}

\subsection{Limitações, ética e pendências}
\begin{itemize}[leftmargin=*]
  \item Limite de resolução da modularidade~\cite{fortunato2007}. \pendente{discutir}
  \item LGPD e anonimização dos dados. \pendente{discutir}
\end{itemize}

\subsection{Registro}
\begin{description}[leftmargin=2.6cm,style=nextline]
  \registro{AAAA-MM-DD}{\pendente{primeira entrada}}
\end{description}

% =====================================================================
% Frente 5 --- Influenciadores (Ravi)
% =====================================================================
\section{Frente 5 --- Influenciadores}

\noindent\textbf{Responsável:} Ravi.
\textbf{No relatório:} problema aplicado e resultados dos influenciadores; referências em BibTeX.
\textbf{Na Etapa 1:} descrição formal do problema aplicado.

\subsection{O que foi implementado}

{\small\begin{tabularx}{\linewidth}{@{}L{3.9cm} L{4.7cm} Y L{1.9cm}@{}}
\toprule
Função & Arquivo & Complexidade & Situação \\
\midrule
\codigo{limiar\_epidemico} & \codigo{aplicacao/influencia.py} & $O(n)$ & \afazer \\
\codigo{simular\_sir} & \codigo{aplicacao/influencia.py} & \falta & \afazer \\
\codigo{poder\_de\_espalhamento} & \codigo{aplicacao/influencia.py} & \falta & \afazer \\
\codigo{tau\_de\_kendall} (tau-b) & \codigo{aplicacao/influencia.py} & $O(n^2)$ ingênuo~\cite{kendall1938} & \afazer \\
\codigo{comparar\_rankings} & \codigo{aplicacao/influencia.py} & --- & \afazer \\
\bottomrule
\end{tabularx}}

\subsection{Problema aplicado}
\pendente{a pergunta, a entrada, a saída e como validar --- texto da descrição formal da Etapa 1}

\subsection{Decisões e justificativas}
\begin{itemize}[leftmargin=*]
  \item SIR em tempo discreto, com $\beta$ um pouco acima do limiar epidêmico (0,0079 na Twitch PT),
        por exemplo 0,012, e $\gamma = 1$, seguindo~\cite{kitsak2010}. \pendente{confirmar os valores usados}
\end{itemize}

\subsection{Corretude: testes e resultados}

{\small\begin{tabularx}{\linewidth}{@{}Y L{2.5cm} L{2.2cm} L{3.8cm}@{}}
\toprule
Teste & Base & Obtido & Referência \\
\midrule
SIR com $\beta = 0$ atinge só os iniciais; $\beta = \gamma = 1$ num conexo atinge todos & sintético & \falta & --- \\
Kendall de rankings iguais e invertidos & sintético & \falta & 1 e $-1$ \\
Limiar epidêmico & Twitch PT & \falta & 0,0079 \\
$\tau$ entre grau e visualizações & Twitch PT & \falta & $\approx$ 0,43 (NetworkX) \\
$\tau$ entre PageRank / k-core e visualizações & Twitch PT & \falta & $\approx$ 0,44 / 0,40 \\
\bottomrule
\end{tabularx}}

\subsection{Resultados}
\pendente{tabela com o $\tau$ de cada centralidade contra o alcance SIR e contra as visualizações;
o k-core prevê os espalhadores melhor que o grau?}

\subsection{Limitações e pendências}
\begin{itemize}[leftmargin=*]
  \item \pendente{ex.: custo do SIR por Monte Carlo; número de rodadas usado}
\end{itemize}

\subsection{Registro}
\begin{description}[leftmargin=2.6cm,style=nextline]
  \registro{AAAA-MM-DD}{\pendente{primeira entrada}}
\end{description}

% =====================================================================
% Frente 6 --- Interface e visualização (Kaio)
% =====================================================================
\section{Frente 6 --- Interface e visualização}

\noindent\textbf{Responsável:} Kaio.
\textbf{No relatório:} interface; conclusão.
\textbf{Na Etapa 1:} diagrama de casos de uso e esboço das telas.

\subsection{O que foi implementado}

{\small\begin{tabularx}{\linewidth}{@{}L{3.9cm} L{4.7cm} Y L{1.9cm}@{}}
\toprule
Item & Arquivo & Complexidade / observação & Situação \\
\midrule
Diagrama de casos de uso & \codigo{documentos/uml/} & 17 casos de uso & \andamento \\
\codigo{fruchterman\_reingold} & \codigo{visualizacao/}\newline\codigo{posicionamento.py} & $O(n^2)$ por iteração~\cite{fruchterman1991} & \afazer \\
Tela: visão geral & \codigo{interface/} & --- & \afazer \\
Tela: rede & \codigo{interface/} & --- & \afazer \\
Tela: centralidades & \codigo{interface/} & --- & \afazer \\
Tela: comunidades & \codigo{interface/} & --- & \afazer \\
Tela: influenciadores & \codigo{interface/} & --- & \afazer \\
Tela: desempenho & \codigo{interface/} & --- & \afazer \\
\bottomrule
\end{tabularx}}

\subsection{Decisões e justificativas}
\begin{itemize}[leftmargin=*]
  \item Streamlit para as telas e Plotly para os gráficos; a interface só exibe resultados
        calculados pelo nosso código. \pendente{completar}
  \item \pendente{ex.: \codigo{st.cache\_resource} para grafos e \codigo{st.cache\_data} para resultados}
\end{itemize}

\subsection{Corretude: testes e resultados}

{\small\begin{tabularx}{\linewidth}{@{}Y L{2.5cm} L{2.2cm} L{3.8cm}@{}}
\toprule
Teste & Base & Obtido & Referência \\
\midrule
Posição para todos os nós, sem valores inválidos & sintético & \falta & --- \\
Mesma \codigo{semente} repete o layout & sintético & \falta & --- \\
\bottomrule
\end{tabularx}}

\subsection{Figuras e tabelas para o relatório}
\begin{itemize}[leftmargin=*]
  \item Diagrama de casos de uso (\codigo{documentos/uml/casos\_de\_uso.jpg}).
  \item \pendente{prints das 6 telas}
\end{itemize}

\subsection{Limitações e pendências}
\begin{itemize}[leftmargin=*]
  \item \pendente{ex.: redes grandes desenhadas só com top-$k$ ou k-core}
\end{itemize}

\subsection{Registro}
\begin{description}[leftmargin=2.6cm,style=nextline]
  \registro{2026-11-10}{Diagrama de casos de uso em PlantUML e descrição dos casos de uso (PR \#5, em revisão).}
\end{description}


\bibliographystyle{plain}
\bibliography{referencias}

\end{document}
